\documentclass[10pt,a4paper]{amsart}
\usepackage[margin=19mm]{geometry}
\usepackage{amsmath,amssymb,mathtools}
\usepackage[scr=rsfs,cal=euler]{mathalfa}
\usepackage{xfrac}
\usepackage[unicode,hidelinks,bookmarks=false]{hyperref}
\newcommand{\ie}{i.e.\ }
\newcommand{\cf}{cf.\ }
\newcommand{\resp}{respectively\ }
\newcommand{\Subsection}{\subsection}
\providecommand{\nobreakdash}{\nobreak\hbox{-}}
\setlength{\emergencystretch}{2em}
\multlinegap0pt
\usepackage{amssymb}
\usepackage{mathtools}
\usepackage{bm}
\usepackage{tikz}
\usepackage{algorithm}\usepackage{algpseudocode}
\newtheorem{lemma}{Lemma}
\title{Causal Architecture in Hidden Quantum Markov Models}
\author{Abdessatar Souissi, Abdessatar Barhoumi}
\date{Source: arXiv:2602.19120v2 (CC BY 4.0)}
\begin{document}
\maketitle

\noindent\emph{Source and licence.} Causal Architecture in Hidden Quantum Markov Models. Version arXiv:2602.19120v2 (CC BY 4.0), distributed by arXiv under the Creative Commons Attribution 4.0 International licence, http://creativecommons.org/licenses/by/4.0/, as recorded in the arXiv licence metadata for this exact version and shown on the abstract page https://arxiv.org/abs/2602.19120v2. Full licence notice: \texttt{licenses/arxiv-2602.19120-CC-BY-4.0.txt}.\par\smallskip

\noindent\textit{Editorial scope.} Counted unit: the article's lemma lem:dual-block-maps (source0.tex lines 471--530). The article's complete original proof of it is delivered with it (source0.tex lines 532--630, one \texttt{proof} environment of 99 lines). No mathematics is written, repaired or re-derived: the statement and the proof are the article's own lines, in the article's own notation, and no retained line is substituted or re-flowed.  Retained uncounted as the source-local context the proof needs: lines 157--159; lines 303--307; lines 309--313.  Nothing was omitted from any retained span, so the package carries no omission record.  No retained line contains a citation, so the wrapper carries no bibliography and no bibliography entry.  No retained line contains a figure environment, an image-inclusion command or drawing code, so no image is delivered and the package declares no asset dependency.  Editorial formatting only, in addition: an AMS \texttt{amsart} preamble that restates the article's own declarations and macros used by the retained lines, namely the environments \texttt{lemma} on separate counters, together with the standard packages those lines need.  The retained environments print in this document as Lemma 1 in order of appearance, whereas the article prints the same items with the numbering of its own full document.  The complete native source of the article as distributed in the arXiv e-print for this exact version is bundled unchanged as \texttt{sources/195086/source0.tex}.
\begin{equation}\label{eq_HS}
\operatorname{Tr}(\mathcal{E}_*(\rho)X)=\operatorname{Tr}(\rho\,\mathcal{E}(X))
\end{equation}

\begin{equation}\label{eq_F}
\mathcal{F}_{a_n,b_n}^{(n)}(a_{n+1})
:=
\mathcal{E}_{H;n}\bigl(\mathcal{E}_{H,O;n}(a_n\otimes b_n)\otimes a_{n+1}\bigr),
\end{equation}

\begin{equation}\label{eq_G}
\mathcal{G}_{a_n,b_n}^{(n)}(a_{n+1})
:=
{\mathcal{E}}_{H,O;n}\bigl(\mathcal{E}_{H;n}(a_n\otimes a_{n+1})\otimes b_n\bigr),
\end{equation}

\begin{lemma}\label{lem:dual-block-maps} 
Fix a time step $n$ and suppose that the hidden transition and emission expectations admit minimal Kraus decompositions
\[
\mathcal{E}_{H;n}(X)
=
\sum_{\alpha=1}^{r_H} K_{H;\alpha}^{*} X K_{H;\alpha},
\qquad
\mathcal{E}_{H,O;n}(Y)
=
\sum_{\beta=1}^{r_O} K_{H,O;\beta}^{*} Y K_{H,O;\beta},
\]
where $K_{H;\alpha}:\mathcal{H}_n\to\mathcal{H}_n\otimes\mathcal{H}_{n+1}$ and $K_{H,O;\beta}:\mathcal{H}_n\to\mathcal{H}_n\otimes\mathcal{K}_n$ satisfy
\(
\sum_{\alpha}K_{H;\alpha}^{*}K_{H;\alpha}
=
\sum_{\beta}K_{H,O;\beta}^{*}K_{H,O;\beta}
=
\mathbb{I}_{\mathcal{H}_n}
\).
Let $\mathcal{F}_{a,b}^{(n)}$ and $\mathcal{G}_{a,b}^{(n)}$ be the block maps defined in \eqref{eq_F} and \eqref{eq_G}, and let
\[
\mathcal{F}_{a,b}^{(n)*},\mathcal{G}_{a,b}^{(n)*}:\mathfrak{S}_{\le 1}(\mathcal{H}_n)\to\mathfrak{S}(\mathcal{H}_{n+1})
\]
denote their Schr\"odinger-picture duals with respect to the Hilbert–Schmidt pairing \eqref{eq_HS}. Then, for every $\rho\in\mathfrak{S}(\mathcal{H}_n)$,
\begin{align} \label{Fab*}
\mathcal{F}_{a,b}^{(n)*}(\rho)
&=
\sum_{\alpha=1}^{r_H}
\operatorname{Tr}_{\mathcal{H}_n}\Bigl[
K_{H;\alpha}\,\rho\,K_{H;\alpha}^{*}
\bigl(\mathcal{E}_{H,O;n}(a\otimes b)\otimes\mathbb{I}_{\mathcal{H}_{n+1}}\bigr)
\Bigr]\nonumber\\[0.4em]
&=
\sum_{\alpha=1}^{r_H}\sum_{\beta=1}^{r_O}
\operatorname{Tr}_{\mathcal{H}_n}\Bigl[
K_{H;\alpha}\,\rho\,K_{H;\alpha}^{*}
\Bigl(
 K_{H,O;\beta}^{*}(a\otimes b)K_{H,O;\beta}
\otimes\mathbb{I}_{\mathcal{H}_{n+1}}
\Bigr)
\Bigr]\\[0.4em]
\label{Gab*}\mathcal{G}_{a,b}^{(n)*}(\rho)
&=
\sum_{\beta=1}^{r_O}
\operatorname{Tr}_{\mathcal{H}_n}\Bigl[
K_{H,O;\beta}\,\rho\,K_{H,O;\beta}^{*}
\bigl(\mathcal{E}_{H;n}(a\otimes\mathbb{I}_{\mathcal{H}_{n+1}})\otimes b\bigr)
\Bigr]\nonumber\\
&=
\sum_{\alpha=1}^{r_H}\sum_{\beta=1}^{r_O}
\operatorname{Tr}_{\mathcal{H}_n}\Bigl[
K_{H,O;\beta}\,\rho\,K_{H,O;\beta}^{*}
\Bigl(
K_{H;\alpha}^{*}(a\otimes\mathbb{I}_{\mathcal{H}_{n+1}})K_{H;\alpha}
\otimes b
\Bigr)
\Bigr]
\end{align}
where $\operatorname{Tr}_{\mathcal{H}_n}$ denotes the partial trace over the hidden space at time $n$.
\end{lemma}

\begin{proof}
We prove the formula for $\mathcal{F}_{a,b}^{(n)*}$; the argument for $\mathcal{G}_{a,b}^{(n)*}$ is analogous. The duality relation \eqref{eq_HS} states that for all $\rho\in\mathfrak{S}(\mathcal{H}_n)$ and $X\in\mathcal{B}(\mathcal{H}_{n+1})$,
\[
\operatorname{Tr}\bigl(\mathcal{F}_{a,b}^{(n)*}(\rho)\,X\bigr)
=
\operatorname{Tr}\bigl(\rho\,\mathcal{F}_{a,b}^{(n)}(X)\bigr).
\]
Using \eqref{eq_F} and the Kraus representation of $\mathcal{E}_{H;n}$, we obtain
\[
\mathcal{F}_{a,b}^{(n)}(X)
=
\mathcal{E}_{H;n}\bigl(\mathcal{E}_{H,O;n}(a\otimes b)\otimes X\bigr)
=
\sum_{\alpha=1}^{r_H}
K_{H;\alpha}^{*}\bigl(\mathcal{E}_{H,O;n}(a\otimes b)\otimes X\bigr)K_{H;\alpha}.
\]
Hence
\[
\operatorname{Tr}\bigl(\rho\,\mathcal{F}_{a,b}^{(n)}(X)\bigr)
=
\sum_{\alpha=1}^{r_H}
\operatorname{Tr}\Bigl(\rho\,K_{H;\alpha}^{*}\bigl(\mathcal{E}_{H,O;n}(a\otimes b)\otimes X\bigr)K_{H;\alpha}\Bigr).
\]
By cyclicity of the trace, each summand can be rewritten as
\[
\operatorname{Tr}\Bigl(K_{H;\alpha}\rho K_{H;\alpha}^{*}\bigl(\mathcal{E}_{H,O;n}(a\otimes b)\otimes X\bigr)\Bigr).
\]
We now interpret this as a trace over $\mathcal{H}_n\otimes\mathcal{H}_{n+1}$ and factor it via the partial trace. For any operator $M$ on $\mathcal{H}_n\otimes\mathcal{H}_{n+1}$ and $X$ on $\mathcal{H}_{n+1}$, one has
\[
\operatorname{Tr}_{\mathcal{H}_n\otimes\mathcal{H}_{n+1}}\bigl[M(\mathbb{I}_{\mathcal{H}_n}\otimes X)\bigr]
=
\operatorname{Tr}_{\mathcal{H}_{n+1}}\Bigl[\operatorname{Tr}_{\mathcal{H}_n}(M)\,X\Bigr],
\]
so with
\[
M
=
K_{H;\alpha}\rho K_{H;\alpha}^{*}
\bigl(\mathcal{E}_{H,O;n}(a\otimes b)\otimes\mathbb{I}_{\mathcal{H}_{n+1}}\bigr)
\]
we obtain
\[
\operatorname{Tr}\Bigl(K_{H;\alpha}\rho K_{H;\alpha}^{*}\bigl(\mathcal{E}_{H,O;n}(a\otimes b)\otimes X\bigr)\Bigr)
=
\operatorname{Tr}_{\mathcal{H}_{n+1}}\Bigl[
\operatorname{Tr}_{\mathcal{H}_n}\Bigl(
K_{H;\alpha}\rho K_{H;\alpha}^{*}
\bigl(\mathcal{E}_{H,O;n}(a\otimes b)\otimes\mathbb{I}_{\mathcal{H}_{n+1}}\bigr)
\Bigr)X
\Bigr].
\]
Summing over $\alpha$ and using linearity of the partial trace, we arrive at
\[
\operatorname{Tr}\bigl(\rho\,\mathcal{F}_{a,b}^{(n)}(X)\bigr)
=
\operatorname{Tr}_{\mathcal{H}_{n+1}}\Bigl(
\Bigl[
\sum_{\alpha=1}^{r_H}
\operatorname{Tr}_{\mathcal{H}_n}\Bigl(
K_{H;\alpha}\rho K_{H;\alpha}^{*}
\bigl(\mathcal{E}_{H,O;n}(a\otimes b)\otimes\mathbb{I}_{\mathcal{H}_{n+1}}\bigr)
\Bigr)
\Bigr]X
\Bigr).
\]
By the defining duality, the square bracket must coincide with $\mathcal{F}_{a,b}^{(n)*}(\rho)$, since this identity holds for all $X$. Thus
\[
\mathcal{F}_{a,b}^{(n)*}(\rho)
=
\sum_{\alpha=1}^{r_H}
\operatorname{Tr}_{\mathcal{H}_n}\Bigl[
K_{H;\alpha}\rho K_{H;\alpha}^{*}
\bigl(\mathcal{E}_{H,O;n}(a\otimes b)\otimes\mathbb{I}_{\mathcal{H}_{n+1}}\bigr)
\Bigr].
\]
Finally, inserting the Kraus expansion of $\mathcal{E}_{H,O;n}$,
\[
\mathcal{E}_{H,O;n}(a\otimes b)
=
\sum_{\beta=1}^{r_O}K_{H,O;\beta}^{*}(a\otimes b)K_{H,O;\beta},
\]
we obtain the first formula stated in the lemma.

The derivation for $\mathcal{G}_{a,b}^{(n)*}$ is identical in structure. Starting from \eqref{eq_G}, using the Kraus decomposition of $\mathcal{E}_{H,O;n}$, cycling the trace, and factorising via partial traces leads to
\[
\operatorname{Tr}\bigl(\rho\,\mathcal{G}_{a,b}^{(n)}(X)\bigr)
=
\operatorname{Tr}_{\mathcal{H}_{n+1}}\Bigl(
\Bigl[
\sum_{\beta=1}^{r_O}
\operatorname{Tr}_{\mathcal{H}_n}\Bigl(
K_{H,O;\beta}\rho K_{H,O;\beta}^{*}
\bigl(\mathcal{E}_{H;n}(a\otimes\mathbb{I}_{\mathcal{H}_{n+1}})\otimes b\bigr)
\Bigr)
\Bigr]X
\Bigr)
\]
for all $X$. By uniqueness of the Hilbert–Schmidt dual, the bracketed operator is exactly $\mathcal{G}_{a,b}^{(n)*}(\rho)$, and expanding $\mathcal{E}_{H;n}$ in its Kraus form yields the second formula. This completes the proof.
\end{proof}

\end{document}
